\section{Introduction}\label{sec:intro}

\subsection{Problem, and what the companion papers settle}\label{sec:intro:problem}

For a linear system $(A,C)$ with a prescribed block $L_0$ of functional rows, a design is a matrix
$\mathcal L=[L_0;R]$ with $q$ independent rows, and it is feasible when
\begin{equation}\label{eq:5a}
\rank\begin{bmatrix}\mathcal LA\\ \mathcal L\\ CA\\ C\end{bmatrix}
=\rank\begin{bmatrix}\mathcal L\\ CA\\ C\end{bmatrix},
\end{equation}
\begin{equation}\label{eq:5b}
\rank\begin{bmatrix}\lambda\mathcal L-\mathcal LA\\ CA\\ C\end{bmatrix}
=\rank\begin{bmatrix}\mathcal L\\ CA\\ C\end{bmatrix}\ \ \forall\lambda\in\mathbb C .
\end{equation}
The companion papers settle the questions one usually asks about \eqref{eq:5a}--\eqref{eq:5b}. The
first \cite{companion1} decides feasibility for a prescribed order, after compressing the free rows to
a subspace of a compressed row-pair space; the second \cite{companion2} computes the minimum order and
constructs a minimum-order observer from a subspace $W$ of that space, with the structural bound
$\omega+1\le\dim W\le 2\omega$.

Neither asks what the designs themselves look like. This paper describes them for the smallest layer
in which a direction exists, that is, when the kernel of the compressed pencil is one-dimensional. We
call the set of all admissible designs of that layer its \emph{algebraic parameter space}, and we
describe it as a finite union of fibre bundles; each bundle sits over a \emph{locally closed algebraic
stratum} of a projective eigenspace, and its fibres are Grassmannians. No single closed subvariety
describes the whole parameter space, and the non-eigen stratum in particular is open rather than
closed, so the stratum-by-stratum statement is the natural one.

\subsection{Contributions}\label{sec:intro:contributions}

\begin{enumerate}[label=(\Alph*),leftmargin=*]
\item \emph{The kernel direction is an eigenvector.} The kernel line of $H=\mathcal LN$ is not free: it
is generated by a real eigenvector $v=Na$ of the orthogonal compression $M=\pi_{K_0}A|_{K_0}$, where
$a\in\mathbb R^k$ collects the coordinates of the direction in the basis $N$ of $K_0$. With prescribed
rows the admissible directions are further restricted to
$E_\mu^0=E_\mu\cap\ker(L_0N)\cap\ker(L_0AN)$ (Theorem~\ref{sec3:thm:direction},
Lemma~\ref{sec3:lem:prescribed}). This turns "choose $q$ rows" into "choose a direction, then a
subspace over it".
\item \emph{The algebraic parameter space is a finite union of relative Grassmann bundles.} Over each
locally closed stratum of each projective eigenspace --- the eigen directions and the non-eigen
directions --- the admissible designs form a relative Grassmann bundle whose fibre dimension and
nonemptiness condition are explicit (Theorem~\ref{sec3:thm:param}). The fibre dimension is
$(s-\delta)(r(v)-s)$, where $s$ is the number of rows beyond the compulsory $k-1$, $r(v)\in\{\rho,\rho-1\}$
is the rank of the quotient bundle, and $\delta$ measures the prescribed rows modulo $K_0$.
\item \emph{One computation per chart, valid at every point of it.} On a chart of the parameter space,
the maximal minors have no common root over the function field if and only if \emph{every} point of
that chart is feasible. The step from a generic point to the whole chart is a full-rank
leading-coefficient lemma (Lemma~\ref{sec4:lem:fullrank}); without it a generic test certifies only a
generic point. This is what makes the chartwise computation a decision procedure rather than a
sampling heuristic, and it is the property that the counterexample below would otherwise violate.
\item \emph{Fixed-$p$ tractability, and where it stops.} For fixed output dimension $p$ and a
geometrically simple real spectrum of $M$, the decision is polynomial in the input bit size
(Theorem~\ref{sec5:thm:poly}), with the exponent depending on $p^2$. The hypothesis governs the
dimension of the direction parameter, and it cannot simply be dropped; no hardness statement for the
repeated case is claimed.
\item \emph{A certified counterexample.} Inside a two-dimensional eigenspace one direction is feasible
and five others are not, with complete rational data and exact certificates
(Section~\ref{sec:counterexample}). Hence "one eigenvector per eigenvalue" is not a valid reduction.
\item \emph{A duality with the compressed formulation.} In this layer the projective direction is the
annihilator of the common projection used in \cite{companion2}: with $U=\pi_1(W)=\pi_2(W)$ one has
$U=a^\circ$ and $[a]=\mathbb P(\Ann U)$ (Proposition~\ref{sec3:prop:duality}).
\end{enumerate}

\subsection{Strength of the claims}\label{sec:intro:claims}

Three statements of different strength are made, and they are kept apart throughout: (i) an exact
decision procedure for the whole layer, in exact arithmetic and with no complexity claim
(Section~\ref{sec4}); (ii) polynomial-time decidability on the geometrically simple subclass, as a
promise-class result whose algorithm reports that an input is outside its scope rather than declaring
infeasibility (Section~\ref{sec5}); (iii) a counterexample refuting one specific reduction
(Section~\ref{sec:counterexample}). Nothing is claimed about the complexity of the repeated case.

\subsection{Notation}\label{sec:intro:notation}

The standard Euclidean inner product on $\mathbb R^n$ is used throughout, so restrictions to subspaces
are nondegenerate; inputs are rational and $\mathcal L$ has independent rows. Table~\ref{tab:notation}
collects the symbols; it is repeated nowhere else and is used as the reference for the whole paper.
Two points of the table deserve emphasis. First, $v$ is an ambient vector in $K_0\subseteq\mathbb R^n$
while $a\in\mathbb R^k$ is its coordinate vector, $v=Na$; the kernel direction is a line in the first
space and a projective point in the second, and it is the second that appears in the duality of
contribution~(F). Second, $\rho=\rank[C;CA]$ is the rank of the compressed data and is not the layer
index, which is $c=k-\omega$.

\begin{table*}[t]
\centering
\caption{Symbols used throughout, in two blocks read row by row. The third cell of each block gives
the space and its dimension; ``---'' marks a scalar or an index.}
\label{tab:notation}
\small
\begin{tabular}{@{}llllll@{}}
\toprule
Symbol & Object & Space, dimension & Symbol & Object & Space, dimension \\
\midrule
$K_0$ & $\ker[C;CA]$, unobservable subspace & $\mathbb R^n$, $k$ &
$E_\mu$ & $\ker(M-\mu I)$, $\mu$ real & $K_0$, $d_\mu$ \\
$k$ & $\dim K_0$ & --- &
$E_\mu^0$ & $E_\mu\cap\ker(L_0N)\cap\ker(L_0AN)$ & $K_0$, $d_\mu^0$ \\
$\rho$ & $\rank[C;CA]=n-k$ & --- &
$v$, $a$ & direction and coordinates, $v=Na$ & $K_0$ resp.\ $\mathbb R^k$ \\
$N$ & basis matrix of $K_0$ & $\mathbb R^{n\times k}$ &
$U_v$ & $\operatorname{span}\{v,Av\}$ & $\mathbb R^n$, $1$ or $2$ \\
$H$, $H_A$ & $\mathcal LN$, $\mathcal LAN$ & $\mathbb R^{q\times k}$ &
$D_v$ & $K_0\cap U_v^\perp$ & $K_0$, $k-1$ \\
$\omega$ & $\rank H=\dim(V\cap K_0)$ & --- &
$\mathcal Q_v$ & $U_v^\perp/D_v$ & quotient, $r(v)$ \\
$c$ & layer index $k-\omega$ & --- &
$V_0$ & $\row L_0$, prescribed rows & $\mathbb R^n$, $r_0$ \\
$M$ & $\pi_{K_0}A|_{K_0}$ & $\operatorname{End}K_0$, $k\times k$ &
$\delta$ & $\dim V_0-\dim(V_0\cap K_0)$ & --- \\
$A_{21}$ & $\pi_{K_0^\perp}A|_{K_0}$ & $K_0\to K_0^\perp$ &
$s$ & $q-k+1$, rows beyond $k-1$ & --- \\
$r(v)$ & $\rho$ (eigen), $\rho-1$ (non-eigen) & --- &
$\widetilde r$ & $r(v)-\delta$, rank left for free rows & --- \\
\bottomrule
\end{tabular}
\end{table*}
